
\section{Introduction}

Particle-based methods provide a natural framework for modeling systems
composed of discrete interacting entities. Such methods are widely used
to study solids, fluids, granular materials, multiphase flows, molecular
systems, and a broad range of engineering and scientific applications
\cite{Meakin2009,onate2011,allen2017}. By explicitly representing
individual interacting elements, particle methods provide a flexible
means of studying complex collective behavior that emerges from local
particle interactions.

As simulation fidelity increases, the number of particles required to
accurately represent physical systems may grow into the millions or
billions. Higher particle counts enable finer spatial resolution and
improved physical fidelity. However, these benefits come at the cost of
substantially increased computational requirements. One of the most
significant of these costs is collision detection.

A naïve all-pairs collision search scales as \(O(N^2)\), making direct
evaluation impractical for large particle counts. Broad-phase
collision-detection algorithms reduce this cost by restricting collision
testing to nearby particles and therefore play a central role in the
scalability of large-scale particle simulations \cite{ericson2004}.

One approach to improving collision-detection performance has been the
migration of broad-phase and narrow-phase processing to graphics
processing units (GPUs). Modern GPUs provide substantial computational
throughput and memory bandwidth and are widely used for large-scale
particle simulation. Consequently, numerous GPU-based collision-
detection methods have been developed to accelerate candidate
generation and collision evaluation. Despite these advances,
occupancy generation remains a dedicated collision-detection operation
that is largely independent of the rendering pipeline.

The approach proposed in this work addresses this challenge through a
new execution architecture. Rather than treating occupancy generation
as a dedicated simulation operation, the proposed graphics-based
Parallel Collision Detection (gPCD) framework integrates broad-phase
occupancy generation directly into the graphics pipeline, allowing
occupancy construction and rendering to proceed concurrently while
producing the potentially colliding set required for narrow-phase
evaluation. This organization enables collision-detection work to be performed using
graphics-pipeline resources traditionally dedicated to rendering.


The principal contributions of this work are:

\begin{enumerate}

\item A graphics-pipeline occupancy-generation framework that integrates
broad-phase construction directly into the rendering pipeline.

\item A unified graphics-compute collision-detection architecture that
constructs the potentially colliding set during graphics execution and
performs narrow-phase evaluation using compute shaders.

\item A complexity analysis of occupancy generation, narrow-phase
evaluation, and storage requirements.

\item Experimental validation demonstrating near-linear scaling
behavior across graphics, compute, and total execution times.

\item Demonstration of multi-million-particle collision detection on a
single CPU--GPU system.

\end{enumerate}
